\section{Introduction}
\label{sec:introduction}

Wikipedia is a widely used source of knowledge, attracting billions of monthly visitors~\citep{statista_wikipedia_traffic_2024}. Although initially criticized for reliability, the English Wikipedia later gained broad acceptance as a reputable source~\citep{economist_wikipedia_2021}. Beyond public use, it plays a central role in natural language processing (NLP) research: Wikipedia is used to train large language models (LLMs), provide ground truth for retrieval-augmented generation (RAG) systems~\citep{semnani-etal-2023-wikichat, lewis2020retrieval, guu2020realm, zhang-etal-2024-spaghetti}, and supply gold answers for question answering and fact verification.

Given this reliance, ensuring Wikipedia's accuracy is critical. We focus specifically on internal inconsistencies: contradictory facts within Wikipedia that indicate errors requiring correction through consultation of original sources. In a crowdsourced repository, inconsistencies can arise from outdated information, limited awareness of related content during editing, or simple human error.

The corpus's vast scale makes comprehensive verification challenging for both humans and automated tools.
While Wikipedia is often used to detect hallucinations in LLMs, we instead leverage LLMs to detect inconsistencies in a human-curated corpus.
Our contributions are as follows:

\textbf{We formalize the task of Corpus-Level Inconsistency Detection (\task).} Given a fact from a corpus, the goal is to identify at least one other fact within the same corpus that contradicts it. While inconsistency detection has been studied at the sentence-pair and document levels, corpus-level detection remains largely unexplored. Recent work examines inconsistencies between retrieved information and the internal knowledge of LLMs~\citep{su2024textttconflictbank, jin2024tugofwar, xie2023adaptive, wang2025retrievalaugmented}, but often relies on synthetic edits or focuses solely on temporal drift~\citep{marjanovic-etal-2024-dynamicqa}. Our task differs from traditional knowledge-intensive settings~\citep{petroni-etal-2021-kilt}, such as question answering~\citep{kwiatkowski-etal-2019-natural, joshi-etal-2017-triviaqa} and fact verification~\citep{thorne-etal-2018-fever, jiang-etal-2020-hover}, which typically assume corpus consistency: finding a single supporting or refuting evidence item is sufficient. This assumption breaks down when the corpus itself contains contradictions. Figure~\ref{fig:task} illustrates this distinction using an example from the \feverous dataset~\cite{aly-etal-2021-fact} and contrasts it with how \system addresses the same case.

\textbf{We propose \system (Corpus-Level Assistant for Inconsistency REcognition), a system for surfacing inconsistencies in large corpora.} To support non-expert users, \system finds and displays not only candidate contradictions but also disambiguating context and explanations of specialized terminology. It features an interactive interface implemented as a browser extension that surfaces potential inconsistencies to Wikipedia visitors. In a user study with eight experienced editors, participants identified 64.7\% more inconsistencies within the same amount of time when using \system than when using search engines.

\textbf{We provide the first lower bound on the inconsistency rate in the English Wikipedia and in two widely used Wikipedia-based NLP benchmarks.} Through manual verification of \system outputs, we estimate that approximately 3.3\% of facts in the English Wikipedia contradict other statements in the corpus. To our knowledge, this is the first systematic attempt to quantify corpus-level inconsistencies in Wikipedia. In AmbigQA~\cite{min-etal-2020-ambigqa}, 4.0\% of questions have answers that contradict other content in the same Wikipedia dump, challenging the dataset's assumption of unambiguous, unique answers. In \feverous, 7.3\% of claims labeled as \supports~are contradicted by other evidence within Wikipedia, undermining the standard assumption of corpus consistency in fact verification.

\textbf{We introduce \dataset, a dataset for Corpus-Level Inconsistency Detection on Wikipedia.} \dataset contains inconsistencies identified in the English Wikipedia. Unlike synthetic datasets, it captures genuine ambiguities and complex factual relationships arising in real-world content. To ensure meaningful coverage, we target Wikipedia's Level 5 Vital Articles,\footnote{\url{https://en.wikipedia.org/wiki/Wikipedia:Vital_articles/Level/5}} which are prioritized for improvement by WikiProject Vital Articles and serve as a centralized watchlist of important entries. Is finding inconsistencies in these pages like finding needles in a haystack? With \system-assisted curation, \dataset comprises 955 facts, 34.7\% of which are inconsistent.

\textbf{We evaluate \system and establish strong baselines on \dataset.} On the \dataset test set, \system achieves an AUROC of 75.1\%, outperforming baselines while leaving substantial headroom for future work.
